\documentclass[14pt,a4paper,oneside]{extbook}
\usepackage[utf8]{vietnam}
\usepackage{amsmath,amssymb,mathrsfs,fancyhdr,enumerate,multirow,makecell}
\usepackage{tkz-euclide,tikz-3dplot,tikz,tkz-tab}
\usetikzlibrary{arrows.meta,calc,intersections,angles,quotes,backgrounds,shapes.geometric,patterns,through,shadows,positioning}
\usepackage{yhmath}
\usepackage[top=0.7cm, bottom=0.7cm, left=1.0cm, right=1.0cm, headheight=16pt, headsep=3mm, footskip=5mm]{geometry}
\usepackage[hidelinks,unicode]{hyperref}
\usepackage{qrcode}
\usepackage[many]{tcolorbox}
\usepackage{fontawesome}

\pagestyle{fancy}
\fancyhf{}
\fancyhead[L]{\footnotesize\bfseries\color{teal!85!black}\faBook~SỔ TAY KIẾN THỨC TOÁN 4}
\fancyhead[R]{\footnotesize\bfseries\color{blue!85!black}Lớp Toán Thầy Lĩnh -- loptoanthaylinh.vn}
\fancyfoot[C]{\footnotesize\bfseries\color{gray!80!black}-- Sổ tay kiến thức: Bài 39. Chia cho số có một chữ số --}
\renewcommand{\headrulewidth}{0.4pt}

\begin{document}

% --- BANNER TIÊU ĐỀ BÀI HỌC ---
\noindent
\begin{tcolorbox}[
  enhanced,
  colback=teal!8!white,
  colframe=teal!80!black,
  arc=3mm, boxrule=1pt,
  left=3mm, right=3mm, top=2mm, bottom=2mm,
  halign=center
]
  {\Large\bfseries\color{teal!90!black}\faBookmark\ BÀI 39: CHIA CHO SỐ CÓ MỘT CHỮ SỐ}\\[1.5pt]
  {\footnotesize\itshape\color{black!75} Quy trình 3 bước chia - nhân - trừ, xử lí chữ số 0 ở thương và điều kiện số dư trên 1 trang A4}
\end{tcolorbox}

\vspace{2mm}

% --- KHỐI 1 & 2: LÝ THUYẾT & KĨ NĂNG ---
\noindent
\begin{minipage}[t]{0.49\textwidth}
  \begin{tcolorbox}[
    enhanced,
    title={\footnotesize\bfseries\color{white}\faSortNumericAsc\ 1. Quy trình chia 3 thao tác},
    colback=blue!2!white,
    colframe=blue!75!black,
    arc=2.5mm, boxrule=0.8pt, top=1.8mm, bottom=1.8mm, left=2mm, right=2mm
  ]
    \footnotesize
    $\bullet$ Thực hiện chia \textbf{từ trái sang phải}:\\
    --- Lượt 1: Lấy 1 chữ số ở số bị chia. Nếu nhỏ hơn số chia thì lấy 2 chữ số.\\
    --- Các lượt sau: Mỗi lần chỉ \textbf{hạ đúng 1 chữ số}.\\
    $\bullet$ Mỗi lượt chia lặp lại đúng 3 thao tác:
    \begin{center}
      \textbf{CHIA} $\rightarrow$ \textbf{NHÂN} $\rightarrow$ \textbf{TRỪ}
    \end{center}
    $\bullet$ Hiệu của phép trừ chính là \textbf{số dư} của lượt chia đó.
  \end{tcolorbox}
\end{minipage}%
\hfill
\begin{minipage}[t]{0.49\textwidth}
  \begin{tcolorbox}[
    enhanced,
    title={\footnotesize\bfseries\color{white}\faShield\ 2. Số dư \& Chữ số 0 ở thương},
    colback=teal!2!white,
    colframe=teal!75!black,
    arc=2.5mm, boxrule=0.8pt, top=1.8mm, bottom=1.8mm, left=2mm, right=2mm
  ]
    \footnotesize
    $\bullet$ \textbf{Điều kiện số dư bất biến}:\\
    \begin{center}
      \textbf{Số dư luôn luôn nhỏ hơn số chia} ($r < b$).
    \end{center}
    --- Số dư lớn nhất có thể có: $\text{Số chia} - 1$.\\
    $\bullet$ \textbf{Chữ số 0 ở thương}: Khi hạ 1 chữ số xuống mà phần bị chia \textbf{nhỏ hơn số chia} thì phải \textbf{viết 0 vào thương} rồi mới được hạ tiếp chữ số sau!
  \end{tcolorbox}
\end{minipage}

\vspace{2mm}

% --- KHỐI 3: SƠ ĐỒ TƯ DUY TIKZ ---
\noindent
\begin{tcolorbox}[
  enhanced,
  title={\footnotesize\bfseries\color{white}\faSitemap\ SƠ ĐỒ TƯ DUY: VÒNG LẶP 3 THAO TÁC CỦA MỖI LƯỢT CHIA},
  colback=yellow!2!white,
  colframe=orange!85!black,
  arc=2.5mm, boxrule=0.8pt, top=1.5mm, bottom=1.5mm, left=1mm, right=1mm,
  halign=center
]
  \begin{tikzpicture}[
    scale=0.9,
    cnode/.style={rectangle, rounded corners=2mm, draw=teal!80!black, fill=teal!10, thick, inner sep=3pt, font=\scriptsize\bfseries, align=center, text width=3.8cm},
    warn/.style={rectangle, rounded corners=2mm, draw=red!80!black, fill=red!10, thick, inner sep=3pt, font=\scriptsize\bfseries, align=center, text width=4.5cm},
    arr/.style={-{Stealth[scale=0.85]}, thick, draw=black!70, shorten >= 2.5pt, shorten <= 2.5pt}
  ]
    \node[cnode, draw=blue!75!black, fill=blue!10] (step1) at (0, 1.0) {\textbf{THAO TÁC 1: CHIA}\\Ước lượng thương\\Lấy chữ số $\ge$ số chia};
    \node[cnode, draw=orange!80!black, fill=orange!10] (step2) at (5.5, 1.0) {\textbf{THAO TÁC 2: NHÂN}\\Thương $\times$ Số chia\\Viết tích dưới phần bị chia};
    \node[cnode, draw=teal!80!black, fill=teal!10] (step3) at (11.0, 1.0) {\textbf{THAO TÁC 3: TRỪ}\\Tìm số dư lượt chia\\Kiểm tra: $\text{Số dư} < \text{Số chia}$};

    \draw[arr] (step1) -- (step2);
    \draw[arr] (step2) -- (step3);

    % Khối cảnh báo và lặp ở dưới
    \node[warn] (w0) at (2.0, -1.0) {\textbf{\color{red!85!black}\faExclamationCircle\ BẪY SỐ 0 Ở THƯƠNG}\\Hạ xuống mà bé hơn số chia\\$\rightarrow$ \textbf{Viết 0 vào thương ngay!}};
    \node[cnode, draw=purple!80!black, fill=purple!10, text width=6.2cm] (next) at (8.5, -1.0) {\textbf{HẠ CHỮ SỐ TIẾP THEO}\\Lặp lại vòng lặp 3 thao tác\\cho đến hết chữ số cuối};

    \draw[arr] (step3.south) -- ++(0, -0.3) -| (next.north);
    \draw[arr] (next.west) -- (w0.east);
    \draw[arr] (w0.north) -- (step1.south);
  \end{tikzpicture}
\end{tcolorbox}

\vspace{2mm}

% --- KHỐI 4: VÍ DỤ MINH HỌA & BẪY TRẮC NGHIỆM ---
\noindent
\begin{minipage}[t]{0.53\textwidth}
  \begin{tcolorbox}[
    enhanced,
    title={\footnotesize\bfseries\color{white}\faCheckCircle\ Ví dụ thực hành chuẩn mực},
    colback=green!2!white,
    colframe=green!60!black,
    arc=2.5mm, boxrule=0.8pt, top=1.8mm, bottom=1.8mm, left=2.2mm, right=2.2mm
  ]
    \footnotesize
    \textbf{Ví dụ}: Đặt tính rồi tính $45\,208 : 4$:\\
    --- $4 : 4 = \mathbf{1}$, $1 \times 4 = 4$, $4 - 4 = 0$.\\
    --- Hạ $5$; $5 : 4 = \mathbf{1}$, $1 \times 4 = 4$, $5 - 4 = 1$.\\
    --- Hạ $2$ được $12$; $12 : 4 = \mathbf{3}$, $3 \times 4 = 12$, $12 - 12 = 0$.\\
    --- Hạ $0$; $0 < 4 \Rightarrow$ \textbf{Viết $\mathbf{0}$ vào thương}!\\
    --- Hạ $8$; $8 : 4 = \mathbf{2}$, $2 \times 4 = 8$, $8 - 8 = 0$.\\
    $\Rightarrow \mathbf{45\,208 : 4 = 11\,302}$ (phép chia hết).
  \end{tcolorbox}
\end{minipage}%
\hfill
\begin{minipage}[t]{0.45\textwidth}
  \begin{tcolorbox}[
    enhanced,
    title={\footnotesize\bfseries\color{white}\faExclamationTriangle\ Mẹo tránh bẫy trắc nghiệm},
    colback=red!2!white,
    colframe=red!70!black,
    arc=2.5mm, boxrule=0.8pt, top=1.8mm, bottom=1.8mm, left=2.2mm, right=2.2mm
  ]
    \footnotesize
    $\bullet$ \textbf{Bẫy quên chữ số 0 ở thương}:\\
    Học sinh hay quên số $0$, ra kết quả $1132$ thay vì $11\,302$ (sai lệch 10 lần!).\\
    $\bullet$ \textbf{Bẫy số dư $\ge$ số chia}:\\
    Nếu sau phép trừ mà số dư $\ge$ số chia chứng tỏ thương ước lượng bị bé, phải tăng thương lên!
  \end{tcolorbox}
\end{minipage}

\end{document}
